\documentclass[11pt,a4paper]{article}

\usepackage[utf8]{inputenc}
\usepackage[T1]{fontenc}
\usepackage{lmodern}
\usepackage[margin=1in,headheight=14pt]{geometry}
\usepackage{microtype}
\usepackage{etoolbox}
\usepackage{xcolor}
\usepackage{xurl}
\usepackage{hyperref}
\usepackage{booktabs}
\usepackage{tabularx}
\usepackage{listings}
\usepackage{tcolorbox}
\usepackage{tikz}
\usepackage{graphicx}
\usepackage{fancyhdr}
\usepackage{lastpage}
\usepackage{titlesec}
\usepackage{enumitem}
\usepackage{setspace}
\usepackage{fontawesome5}

\usetikzlibrary{shapes.geometric, arrows.meta, positioning, calc, backgrounds}

% Line and Paragraph Spacing
\setstretch{1.12}
\setlength{\parskip}{0.5em}
\setlength{\parindent}{0pt}

% List Item Spacing
\setlist[itemize]{itemsep=4pt, topsep=4pt, parsep=2pt, leftmargin=1.5em}
\setlist[enumerate]{itemsep=4pt, topsep=4pt, parsep=2pt, leftmargin=1.5em}

% Custom Color Palette
\definecolor{ghblue}{HTML}{0969DA}
\definecolor{ghdark}{HTML}{1F2328}
\definecolor{ghborder}{HTML}{D0D7DE}
\definecolor{ghcodebg}{HTML}{F6F8FA}
\definecolor{ghgreen}{HTML}{1A7F37}
\definecolor{ghbadgebg}{HTML}{DDF4FF}
\definecolor{ghbadgeborder}{HTML}{54AEFF}

\definecolor{alertred}{HTML}{CF222E}
\definecolor{alertbg}{HTML}{FFEBE9}
\definecolor{alertborder}{HTML}{FF8182}

\definecolor{noteblue}{HTML}{0969DA}
\definecolor{notebg}{HTML}{DDF4FF}
\definecolor{noteborder}{HTML}{54AEFF}

% Hyperlink configuration
\hypersetup{
    colorlinks=true,
    linkcolor=ghblue,
    filecolor=ghblue,
    urlcolor=ghblue,
    citecolor=ghblue,
    pdftitle={Cross-Platform Guide: Build iOS Apps with GitHub and Install Free on iPhone},
    pdfauthor={Mehdi},
    pdfsubject={iOS Development without a Mac},
    pdfkeywords={iOS, Swift, GitHub Actions, Sideloading, Plume Impactor, Sideloadly, Linux}
}

% Page Header & Footer setup (using \pageref* so the last page number is gray, not a blue link)
\pagestyle{fancy}
\fancyhf{}
\renewcommand{\headrulewidth}{0.4pt}
\renewcommand{\footrulewidth}{0.4pt}
\fancyhead[L]{\small\sffamily\color{gray}\faApple\enskip Cross-Platform iOS Build \& Install Guide}
\fancyhead[R]{\small\sffamily\color{gray}\faGithub\enskip GitHub Actions \& Free Sideloading}
\fancyfoot[C]{\small\sffamily\color{gray}Page \thepage\ of \pageref*{LastPage}}

% Plain style for first page
\fancypagestyle{plain}{
    \fancyhf{}
    \renewcommand{\headrulewidth}{0pt}
    \renewcommand{\footrulewidth}{0.4pt}
    \fancyfoot[C]{\small\sffamily\color{gray}Page \thepage\ of \pageref*{LastPage}}
}

% Heading Typography & Spacing
\titleformat{\section}
  {\Large\bfseries\sffamily\color{ghdark}}
  {\thesection}{0.75em}{}
  [\vspace{3pt}\hrule height 0.6pt]

\titleformat{\subsection}
  {\large\bfseries\sffamily\color{ghdark}}
  {\thesubsection}{0.6em}{}

\titleformat{\subsubsection}
  {\normalsize\bfseries\sffamily\color{ghdark}}
  {\thesubsubsection}{0.5em}{}

\titlespacing*{\section}{0pt}{18pt plus 3pt minus 2pt}{9pt plus 2pt minus 2pt}
\titlespacing*{\subsection}{0pt}{14pt plus 3pt minus 2pt}{7pt plus 2pt minus 2pt}
\titlespacing*{\subsubsection}{0pt}{10pt plus 2pt minus 1pt}{5pt plus 1pt minus 1pt}

% Listings configuration with generous spacing
\lstdefinestyle{bashstyle}{
    backgroundcolor=\color{ghcodebg},
    basicstyle=\ttfamily\footnotesize\color{ghdark},
    breakatwhitespace=false,
    breaklines=true,
    columns=fullflexible,
    keepspaces=true,
    showspaces=false,
    showstringspaces=false,
    showtabs=false,
    tabsize=2,
    frame=single,
    rulecolor=\color{ghborder},
    frameround=tttt,
    framexleftmargin=8pt,
    framexrightmargin=8pt,
    framextopmargin=7pt,
    framexbottommargin=7pt,
    aboveskip=10pt,
    belowskip=10pt
}
\lstset{style=bashstyle}

% Prevent code blocks from breaking/dividing across pages
\BeforeBeginEnvironment{lstlisting}{\par\noindent\begin{minipage}{\linewidth}}
\AfterEndEnvironment{lstlisting}{\end{minipage}\par}

% Callout Boxes with generous internal & external spacing (unbreakable so they are never split across pages)
\tcbuselibrary{skins}

\newtcolorbox{notebox}[1][Note]{
    enhanced,
    colback=notebg,
    colframe=noteborder,
    coltitle=noteblue,
    fonttitle=\bfseries\sffamily\small,
    title={\faInfoCircle\enskip #1},
    arc=5pt,
    boxrule=0.8pt,
    left=10pt,
    right=10pt,
    top=8pt,
    bottom=8pt,
    before skip=10pt,
    after skip=10pt
}

\newtcolorbox{importantbox}[1][Important]{
    enhanced,
    colback=alertbg,
    colframe=alertborder,
    coltitle=alertred,
    fonttitle=\bfseries\sffamily\small,
    title={\faExclamationTriangle\enskip #1},
    arc=5pt,
    boxrule=0.8pt,
    left=10pt,
    right=10pt,
    top=8pt,
    bottom=8pt,
    before skip=10pt,
    after skip=10pt
}

\begin{document}

\thispagestyle{plain}

% Title Banner
\begin{center}
    {\huge\bfseries\sffamily Cross-Platform Guide:\\[6pt]
    Build iOS Apps with GitHub \& Install Free on iPhone}\\[12pt]
    {\large\sffamily\color{gray} A complete, beginner-friendly guide to building modern Swift iOS apps in the cloud using \textbf{GitHub Actions} and installing them on physical iPhones for \textbf{free} (no paid Apple Developer account or local Mac required).}
\end{center}

\vspace{10pt}

% Section: Architecture Overview
\phantomsection
\addcontentsline{toc}{section}{Architecture Overview}
\section*{Architecture Overview}

\vspace{4pt}

\begin{center}
\resizebox{\linewidth}{!}{%
\begin{tikzpicture}[
    card/.style={
        rectangle,
        draw=ghborder,
        fill=white,
        thick,
        rounded corners=5pt,
        align=center,
        inner sep=7pt,
        font=\small\sffamily,
        text=ghdark,
        minimum width=6.8cm,
        minimum height=1.12cm
    },
    arrow/.style={
        -{Stealth[length=2.6mm]},
        very thick,
        draw=ghblue
    },
    badge/.style={
        font=\scriptsize\sffamily\bfseries,
        text=ghblue,
        fill=ghbadgebg,
        draw=ghbadgeborder,
        thin,
        rounded corners=3pt,
        inner sep=3pt
    },
    paneltitle/.style={
        font=\bfseries\sffamily\large\color{ghdark},
        align=center
    }
]

% Center coordinates:
% Left Column: x = 0
% Right Column: x = 9.0 cm
% Card width = 6.8 cm
% Panel 1: x in [-3.7, 3.7] (width = 7.4 cm)
% Panel 2: x in [5.3, 12.7] (width = 7.4 cm)
% Corridor: x in [3.7, 5.3], width = 1.6 cm, center x = 4.5 cm

% ==================== PHASE 1: LEFT COLUMN ====================
\node[paneltitle] (t1) at (0, 0.3) {\faCloud\enskip Phase 1: Cloud CI Build};
\node[card] (c1) at (0, -0.85) {\faLaptop\enskip\textbf{Local Workstation (Linux / Windows / macOS)}\\[2pt]\footnotesize Write Swift code using your favorite editor};
\node[card] (c2) at (0, -2.60) {\faGithub\enskip\textbf{GitHub Repository}\\[2pt]\footnotesize Remote source control \& CI workflow triggers};
\node[card] (c3) at (0, -4.35) {\faServer\enskip\textbf{GitHub Actions macOS 15 Runner}\\[2pt]\footnotesize Cloud CI environment executing XcodeGen};
\node[card] (c4) at (0, -6.10) {\faCogs\enskip\textbf{xcodebuild Compilation}\\[2pt]\footnotesize Compiles \& packages unsigned \texttt{MyiOSApp.ipa}};

% Connections - Left Column
\draw[arrow] (c1) -- node[badge] {\faCodeBranch\enskip git push} (c2);
\draw[arrow] (c2) -- node[badge] {\faPlay\enskip Trigger Workflow} (c3);
\draw[arrow] (c3) -- node[badge] {\faHammer\enskip XcodeGen + xcodebuild} (c4);

% ==================== PHASE 2: RIGHT COLUMN ====================
\node[paneltitle] (t2) at (9.0, 0.3) {\faMobile*\enskip Phase 2: Signing \& Device Install};
\node[card] (c5) at (9.0, -0.85) {\faArchive\enskip\textbf{GitHub Artifacts Archive}\\[2pt]\footnotesize Download compiled \texttt{MyiOSApp.ipa} artifact};
\node[card] (c6) at (9.0, -2.60) {\faTools\enskip\textbf{Sideloading Utility}\\[2pt]\footnotesize Plume Impactor (Linux) / Sideloadly (Windows / Mac)};
\node[card] (c7) at (9.0, -4.35) {\faKey\enskip\textbf{Free Apple ID Signing}\\[2pt]\footnotesize Direct 2FA signing via Apple AuthKit servers};
\node[card] (c8) at (9.0, -6.10) {\faMobile*\enskip\textbf{Physical iPhone (iOS 16+)}\\[2pt]\footnotesize App installed \& verified on physical device};

% Connections - Right Column
\draw[arrow] (c5) -- node[badge] {\faDownload\enskip Download \& Extract} (c6);
\draw[arrow] (c6) -- node[badge] {\faShield*\enskip 2FA Apple ID} (c7);
\draw[arrow] (c7) -- node[badge] {\faUsb\enskip Install via USB} (c8);

% ==================== CENTRAL CORRIDOR ROUTING ====================
\draw[arrow] (c4.east) -- (4.5, -6.10) -- node[badge, fill=white, draw=ghbadgeborder, midway] {\faUpload\enskip Upload \texttt{.ipa}} (4.5, -0.85) -- (c5.west);

% ==================== BACKGROUND PANELS ====================
\begin{scope}[on background layer]
    % Left Panel
    \draw[ghborder, fill=ghcodebg, thick, rounded corners=8pt] 
        (-3.7, 0.85) rectangle (3.7, -6.9);
    % Right Panel
    \draw[ghborder, fill=ghcodebg, thick, rounded corners=8pt] 
        (5.3, 0.85) rectangle (12.7, -6.9);
\end{scope}

\end{tikzpicture}%
}
\end{center}

\vspace{8pt}

% Section: Prerequisites
\phantomsection
\addcontentsline{toc}{section}{Prerequisites}
\section*{Prerequisites}

\begin{itemize}
    \item An \textbf{iPhone} running iOS 16 or newer.
    \item A free \textbf{Apple ID} (your standard personal Apple account).
    \item A \textbf{GitHub account}.
    \item A computer running \textbf{Linux}, \textbf{Windows}, or \textbf{macOS}.
    \item A standard USB cable to connect your iPhone to your computer for the initial install.
\end{itemize}

% Section: Step 1
\phantomsection
\addcontentsline{toc}{section}{Step 1: Clone Template, Initialize Git \& Push}
\section*{Step 1: Clone Template, Initialize Git \& Push}

To create your own fresh iOS app using this template, clone this repository, re-initialize a clean git history, and push it to your GitHub account:

\subsection*{1. Clone \& Re-initialize Git}
\begin{lstlisting}[language=bash]
# Clone this starter template
git clone https://github.com/Mehdirben/iosapp.git my-ios-app
cd my-ios-app

# Re-initialize a fresh git history for your project
rm -rf .git
git init
git add .
git commit -m "Initial commit: fresh iOS app template"
\end{lstlisting}

\textit{(Optional: Open \texttt{project.yml} to rename your app and bundle ID to whatever you like)}.

\subsection*{2. Push to Your GitHub}

\begin{notebox}[Actions Quota Considerations]
Public repositories have \textbf{unlimited free GitHub Actions minutes}. Private repositories receive 2,000 free runner minutes/month. Each iOS build takes only $\sim$1.5 minutes.
\end{notebox}

Choose whichever method you prefer:

\subsubsection*{Option A: One-Click GUI via VS Code (Easiest)}
\begin{enumerate}
    \item Open the project directory in VS Code (\texttt{code .}).
    \item Switch to the \textbf{Source Control} tab in the sidebar (\texttt{Ctrl+Shift+G} or \texttt{Cmd+Shift+G}).
    \item Click the \textbf{``Publish to GitHub''} button.
    \item Select \textbf{``Publish to GitHub Public Repository''} or \textbf{``Publish to GitHub Private Repository''}.
    \item VS Code will automatically create the remote repository on your GitHub account and push your code in a single click!
\end{enumerate}

\subsubsection*{Option B: Via Terminal / CLI}
\begin{enumerate}
    \item Create a new repository on \href{https://github.com/new}{GitHub} (Public or Private).
    \item Link your repository and push:
\begin{lstlisting}[language=bash]
git branch -M main
git remote add origin https://github.com/<your-username>/<your-repo-name>.git
git push -u origin main
\end{lstlisting}
\end{enumerate}

Once pushed, GitHub Actions immediately begins compiling your app in the cloud!

% Section: Step 2
\phantomsection
\addcontentsline{toc}{section}{Step 2: Download the Compiled IPA}
\section*{Step 2: Download the Compiled IPA}

\begin{enumerate}
    \item On GitHub, navigate to the \textbf{Actions} tab of your repository.
    \item Select the latest workflow run named \textbf{``Build iOS App (Unsigned IPA)''}.
    \item Once the build finishes (approx. 1.5--2 minutes), scroll down to the \textbf{Artifacts} section.
    \item Click \textbf{\texttt{MyiOSApp-unsigned-ipa}} to download the zip file.
    \item Extract the downloaded zip file on your computer to get \textbf{\texttt{MyiOSApp.ipa}}.
\end{enumerate}

% Section: Step 3
\phantomsection
\addcontentsline{toc}{section}{Step 3: Install the App on Your iPhone}
\section*{Step 3: Install the App on Your iPhone (Choose Your Platform)}

Because the IPA is compiled in CI without Apple certificates, a sideloading utility signs it with your free personal Apple ID before installing it on your device.

% Subsection: Linux
\phantomsection
\addcontentsline{toc}{subsection}{Linux (Fedora, Ubuntu, Debian, Arch)}
\subsection*{\faLinux\enskip Linux (Fedora, Ubuntu, Debian, Arch)}

On Linux, use \href{https://github.com/declaration/impactor}{\textbf{Plume Impactor}} (available directly on \href{https://flathub.org/apps/dev.khcrysalis.PlumeImpactor}{\textbf{Flathub}}) --- a modern, open-source GTK GUI tool that natively supports Apple's AuthKit 2-Factor Authentication (2FA) and free developer signing.

\subsubsection*{1. Install USB Device Communication Tools}
\begin{itemize}
    \item \textbf{Fedora / RHEL}:
\begin{lstlisting}[language=bash]
sudo dnf install -y usbmuxd libimobiledevice libimobiledevice-utils
\end{lstlisting}
    \item \textbf{Ubuntu / Debian}:
\begin{lstlisting}[language=bash]
sudo apt update && sudo apt install -y usbmuxd libimobiledevice6 libimobiledevice-utils
sudo systemctl start usbmuxd
\end{lstlisting}
    \item \textbf{Arch Linux}:
\begin{lstlisting}[language=bash]
sudo pacman -S usbmuxd libimobiledevice
sudo systemctl start usbmuxd
\end{lstlisting}
\end{itemize}

\subsubsection*{2. Pair and Trust Your iPhone}
\begin{enumerate}
    \item Connect your iPhone via USB and unlock the screen.
    \item Pair with Linux using the terminal:
\begin{lstlisting}[language=bash]
idevicepair pair
\end{lstlisting}
    \item A popup saying \textbf{``Trust This Computer?''} will appear on your iPhone screen. Tap \textbf{Trust} and enter your passcode.
    \item Run \texttt{idevicepair pair} once more to confirm:
\begin{lstlisting}[language=bash]
idevicepair pair
\end{lstlisting}
    \textit{(It will output: \texttt{SUCCESS: Paired with device <UDID>})}.
\end{enumerate}

\subsubsection*{3. Install Impactor via Flathub (Recommended)}
Install \textbf{Plume Impactor} directly from Flathub:
\begin{lstlisting}[language=bash]
flatpak install -y flathub dev.khcrysalis.PlumeImpactor

# Grant permission to access local filesystem so you can drag-and-drop .ipa files easily:
flatpak override --user --filesystem=host dev.khcrysalis.PlumeImpactor
\end{lstlisting}
\textit{(Alternative for non-Flatpak systems: download and unpack the AppImage binary from \href{https://github.com/declaration/impactor/releases}{GitHub Releases})}.

\subsubsection*{4. Install the App using Impactor}
\begin{enumerate}
    \item Ensure your iPhone is unlocked and connected via USB.
    \item Launch \textbf{Plume Impactor} from your Applications menu (or run \path{flatpak run dev.khcrysalis.PlumeImpactor}).
    \begin{notebox}[Sequence Note]
    Always unlock and connect your iPhone \textbf{before} opening Impactor so it detects the device on launch.
    \end{notebox}
    \item Your connected iPhone will appear in the device selector.
    \item Select or drag-and-drop your extracted \texttt{MyiOSApp.ipa}.
    \item Sign in with your Apple ID and enter the 6-digit 2FA code sent to your device.
    \item Click \textbf{Install}.
\end{enumerate}

% Subsection: Windows
\phantomsection
\addcontentsline{toc}{subsection}{Windows}
\subsection*{\faWindows\enskip Windows}

On Windows, use \href{https://sideloadly.io/}{\textbf{Sideloadly}} --- a beginner-friendly tool to sign and install apps on your iPhone.

\subsubsection*{1. Why iTunes \& iCloud Are Required on Windows}
Unlike macOS or Linux, Windows has no native drivers for iPhone USB communication, nor libraries to authenticate with Apple's 2FA servers:
\begin{itemize}
    \item \textbf{iTunes} provides Apple's official USB driver (\path{usbaapl64.sys}), allowing your PC to see and talk to the iPhone.
    \item \textbf{iCloud} provides Apple's authentication libraries (\path{ApplePushService.dll}), allowing Sideloadly to perform secure 2-Factor Authentication with Apple.
\end{itemize}

\begin{importantbox}[Do Not Use the Microsoft Store Version]
Microsoft Store apps run inside a locked, isolated sandbox. Sideloadly cannot access the drivers or files inside that sandbox. You \textbf{must} download and install the direct standalone installers from Apple using the links below:
\end{importantbox}

\begin{enumerate}
    \item \textbf{Download \& Install iTunes (64-bit Windows)}: \href{https://www.apple.com/itunes/download/win64}{Direct Apple Download Link}
    \item \textbf{Download \& Install iCloud (Windows)}: \href{https://updates.cdn-apple.com/2020/windows/001-39935-20200911-1A70AA56-F448-11EA-8109-AE43397F938A/iCloudSetup.exe}{Direct Apple Download Link}
\end{enumerate}
\textit{(Restart your computer after installing if prompted).}

\subsubsection*{2. Install Sideloadly}
\begin{enumerate}
    \item Download and install \href{https://sideloadly.io/}{\textbf{Sideloadly (64-bit Windows)}}.
    \item Connect your iPhone to your PC using your USB cable.
    \item Unlock your iPhone screen. When prompted with \textbf{``Trust This Computer?''}, tap \textbf{Trust} and enter your passcode.
\end{enumerate}

\subsubsection*{3. Sign \& Install the App}
\begin{enumerate}
    \item Open \textbf{Sideloadly}. Your connected iPhone will automatically appear in the device dropdown at the top.
    \item Drag and drop your extracted \textbf{\texttt{MyiOSApp.ipa}} into the large app icon box in Sideloadly.
    \item Type your personal Apple ID email address into the \textbf{Apple ID} field.
    \item Click \textbf{Start}.
    \item When prompted, enter your Apple ID password and the 6-digit 2-Factor Authentication code that appears on your iPhone.
    \item Sideloadly will sign the app with your Apple ID and install it directly to your iPhone home screen!
\end{enumerate}

% Subsection: macOS
\phantomsection
\addcontentsline{toc}{subsection}{macOS}
\subsection*{\faApple\enskip macOS}

\begin{enumerate}
    \item Download and install \href{https://sideloadly.io/}{\textbf{Sideloadly}} or \href{https://altstore.io/}{\textbf{AltStore}}.
    \item Connect your iPhone via USB or enable ``Show this iPhone when on Wi-Fi'' in Finder.
    \item Open Sideloadly, drag and drop \texttt{MyiOSApp.ipa}, enter your Apple ID, and click \textbf{Start}.
\end{enumerate}
\textit{(Alternative for macOS: You can also use \href{https://dantheman827.github.io/ios-app-signer/}{iOS App Signer} + Apple Configurator / Xcode).}

% Subsection: SideStore
\phantomsection
\addcontentsline{toc}{subsection}{On-Device / Wire-Free (SideStore)}
\subsection*{\faWifi\enskip On-Device / Wire-Free (SideStore)}

If you want to install and update apps directly on your iPhone without touching a computer every 7 days:

\begin{enumerate}
    \item Perform a one-time install of \href{https://sidestore.io/}{\textbf{SideStore}} using your computer (via Impactor, Sideloadly, or AltServer).
    \item Follow SideStore's on-device setup (installing the WireGuard pairing profile).
    \item From then on:
    \begin{itemize}
        \item Download \texttt{MyiOSApp.ipa} directly from GitHub Actions inside \textbf{Safari on your iPhone}.
        \item Open SideStore, tap \texttt{+}, and select the downloaded \texttt{.ipa}.
        \item SideStore will sign, install, and refresh your apps entirely over local Wi-Fi without needing a PC.
    \end{itemize}
\end{enumerate}

% Section: Step 4
\phantomsection
\addcontentsline{toc}{section}{Step 4: First-Time iOS Device Settings}
\section*{Step 4: First-Time iOS Device Settings}

The first time you launch an app signed with a personal Apple ID, iOS blocks it until you authorize developer access. This is a one-time setup:

\subsection*{1. Trust Your Developer Certificate}
\begin{enumerate}
    \item On your iPhone, open \textbf{Settings} $>$ \textbf{General} $>$ \textbf{VPN \& Device Management}.
    \item Under \textbf{Developer App}, tap your Apple ID.
    \item Tap \textbf{Trust ``[Your Apple ID]''} and confirm.
\end{enumerate}

\subsection*{2. Enable Developer Mode (iOS 16, 17, 18+)}
\begin{enumerate}
    \item Open \textbf{Settings} $>$ \textbf{Privacy \& Security}.
    \item Scroll to the very bottom and tap \textbf{Developer Mode}.
    \item Toggle it \textbf{ON} and tap \textbf{Restart}.
    \item After your iPhone reboots, unlock the screen, tap \textbf{Turn On}, and enter your device passcode.
\end{enumerate}

Your app is now ready to run!

% Section: Troubleshooting
\phantomsection
\addcontentsline{toc}{section}{Troubleshooting}
\section*{Troubleshooting}

\subsection*{Linux: iPhone not detected after restarting laptop or ``Lockdown error -8''}

If your iPhone is not detected after restarting your computer, check these items:

\subsubsection*{1. iPhone is Locked with Passcode (Before First Unlock)}
After rebooting your computer or iPhone, iOS blocks all USB data transfer for security until you unlock the screen.
\begin{notebox}[Rule]
Unlock your iPhone screen with your passcode, then \textbf{unplug and reconnect the USB cable} once so iOS renegotiates the data connection.
\end{notebox}

\subsubsection*{2. Impactor Opened Before Phone Connected}
Plume Impactor inspects USB devices once when it opens. If it was launched before you unlocked/plugged in your iPhone, it won't see it.
\begin{itemize}
    \item \textbf{Fix}: Close Impactor and reopen it after connecting and unlocking your phone.
\end{itemize}

\subsubsection*{3. Root Cause: iPhone USB Hotspot Conflicts with usbmuxd (\texttt{ipheth})}
If Personal Hotspot is enabled (iOS may enable it automatically for trusted computers when cellular data is on), plugging in the iPhone brings up its USB-tethering interface. The Linux \texttt{ipheth} kernel driver claims it and NetworkManager activates it (an \texttt{enp...u5c4i2}-style USB-Ethernet interface appears) --- and while that interface is coming up, it collides with usbmuxd's device handshake: preflight fails with \texttt{lockdown error -8}, the phone starts disconnect/reconnect loops, and usbmuxd stops answering until restarted. This is why the phone is undetected at boot --- or hours after boot whenever it is plugged in with hotspot on --- yet works perfectly with hotspot off.

Confirm it is this case:
\begin{lstlisting}[language=bash]
lsmod | grep ipheth                                  # tethering driver loaded
nmcli device | grep "c4i2"                           # hotspot interface active
journalctl -u usbmuxd -b | grep "lockdown error -8"  # failure signature
\end{lstlisting}

Fixes (pick one):
\begin{itemize}
    \item \textbf{Simplest}: turn Personal Hotspot off (\textbf{Settings} $>$ \textbf{Personal Hotspot}) before connecting the iPhone --- usbmuxd then starts cleanly.
    \item \textbf{Permanent, if you never use USB tethering}: blacklist the \texttt{ipheth} driver so the tethering interface never activates. Wi-Fi hotspot is unaffected:
\begin{lstlisting}[language=bash]
echo "blacklist ipheth" | sudo tee /etc/modprobe.d/blacklist-ipheth.conf
\end{lstlisting}
    Delete that file to re-enable USB tethering later.
    \item \textbf{If you want both}: connect the iPhone with hotspot off (usbmuxd completes its handshake), then enable hotspot --- both coexist fine afterwards.
\end{itemize}

\subsubsection*{4. usbmuxd Wedged Mid-Session (Manual Quick Fix)}
If the phone is attached but invisible (typically after toggling the hotspot while plugged in), restart the daemon:
\begin{lstlisting}[language=bash]
lsusb | grep -i apple   # phone still physically connected?
idevice_id -l           # empty output = usbmuxd lost it
systemctl restart usbmuxd
\end{lstlisting}

Test device communication anytime with:
\begin{lstlisting}[language=bash]
idevice_id -l
\end{lstlisting}
\textit{(Your iPhone's 40-character UDID should print immediately).}

\subsection*{iOS: ``Untrusted Developer'' Error when Opening App}
Go to \textbf{Settings} $>$ \textbf{General} $>$ \textbf{VPN \& Device Management}, select your Apple ID, and tap \textbf{Trust}.

\subsection*{iOS 16+: ``Developer Mode Required''}
Go to \textbf{Settings} $>$ \textbf{Privacy \& Security} $>$ \textbf{Developer Mode}, toggle it \textbf{ON}, and restart your device. Upon reboot, confirm \textbf{Turn On}.

% Section: Apple Rules & Limits
\phantomsection
\addcontentsline{toc}{section}{Apple Free Developer Rules \& Limits}
\section*{Apple Free Developer Rules \& Limits}

\begin{center}
\renewcommand{\arraystretch}{1.35}
\setlength{\tabcolsep}{8pt}
\begin{tabularx}{\textwidth}{l p{4.0cm} X}
\toprule
\textbf{Rule} & \textbf{Limitation} & \textbf{Solution} \\
\midrule
\textbf{Certificate Expiration} & 7 days & Re-sign/reinstall the app once a week using your sideloading tool or SideStore. \\
\addlinespace
\textbf{Max Sideloaded Apps} & 3 active apps simultaneously & If you hit the limit, delete or deactivate an app in your sideloading tool. \\
\addlinespace
\textbf{App ID Registration} & Up to 10 App IDs per 7 days & Keep the \path{bundleIdPrefix} consistent in \path{project.yml}. \\
\bottomrule
\end{tabularx}
\end{center}

% Section: Step 5
\phantomsection
\addcontentsline{toc}{section}{Step 5: Developing \& Adding Features}
\section*{Step 5: Developing \& Adding Features}

You do \textbf{not} need Xcode or a Mac to continue building your application:

\begin{itemize}
    \item \textbf{Edit UI}: Modify \href{file:///home/mehdi/ex/iosapp/Sources/App/ContentView.swift}{\texttt{ContentView.swift}} or create new SwiftUI views in \texttt{Sources/App/}.
    \item \textbf{Add Code Files}: Add any \texttt{.swift} files to \texttt{Sources/App/}. When GitHub Actions runs, \href{https://github.com/yonaskolb/XcodeGen}{\textbf{XcodeGen}} automatically scans the folder and updates the project structure.
    \item \textbf{Deploy Changes}: Commit and push your changes:
\begin{lstlisting}[language=bash]
git commit -am "Add new feature"
git push
\end{lstlisting}
    GitHub Actions will compile the new \texttt{.ipa} automatically. Download and install the update!
\end{itemize}

\end{document}
